\renewcommand{\bibname}{References and further reading}
\begin{thebibliography}{99}
\raggedright
\addcontentsline{toc}{chapter}{References and further reading}
\bibitem{mcs} E. Lehman, F. T. Leighton, and A. R. Meyer, \emph{Mathematics for Computer Science}. The textbook of MIT course 6.1200J, available as a free PDF from MIT OpenCourseWare. A good further source on logic, induction, graphs, matchings, and probability. \url{https://ocw.mit.edu/courses/6-1200j-mathematics-for-computer-science-spring-2024/pages/readings/}
\bibitem{oairepo} OpenAI, \emph{math}: a public repository of mathematical preprints, with Lean formalizations and descriptions of what each formalization covers. \url{https://github.com/openai/math}
\bibitem{oai173} OpenAI, ``Seymour's second-neighborhood conjecture.'' Description of the scope of the Lean formalization, with a link to the accompanying preprint, \emph{A proof of Seymour's second-neighborhood conjecture} (September 2026). \url{https://raw.githubusercontent.com/openai/math/main/lean/docs/173.md}
\bibitem{oai179} OpenAI, ``The circulant Hadamard and Barker-sequence conjectures.'' Description of the scope of the Lean formalization, with a link to the accompanying preprint, \emph{The circulant Hadamard conjecture} (September 2026). \url{https://raw.githubusercontent.com/openai/math/main/lean/docs/179.md}
\bibitem{oaimain} OpenAI, \code{Main.lean}: the Lean file that states and proves the second-neighborhood theorem \code{exists\_goodVertex}. \url{https://raw.githubusercontent.com/openai/math/main/lean/OAI/Combinatorics/SecondNeighborhood/Main.lean}
\bibitem{oaireduction} OpenAI, \code{Reduction.lean}: the Lean file containing the reduction step of the second-neighborhood proof. \url{https://raw.githubusercontent.com/openai/math/main/lean/OAI/Combinatorics/SecondNeighborhood/Reduction.lean}
\bibitem{oaiconfig} OpenAI, configuration file for the formal check of the second-neighborhood theorem, listing the statement file, the proof file, the theorem name, and the permitted axioms. \url{https://raw.githubusercontent.com/openai/math/main/lean/ComparatorChallenges/SeymourSecondNeighborhood.json}
\bibitem{leaninstall} Lean, official installation instructions. \url{https://lean-lang.org/install/}
\bibitem{leanprop} J. Avigad, L. de Moura, S. Kong, and S. Ullrich, \emph{Theorem Proving in Lean 4}, chapter ``Propositions and Proofs.'' \url{https://lean-lang.org/theorem_proving_in_lean4/Propositions-and-Proofs/}
\bibitem{leanquant} J. Avigad, L. de Moura, S. Kong, and S. Ullrich, \emph{Theorem Proving in Lean 4}, chapter ``Quantifiers and Equality.'' \url{https://lean-lang.org/theorem_proving_in_lean4/Quantifiers-and-Equality/}
\bibitem{leantactics} J. Avigad, L. de Moura, S. Kong, and S. Ullrich, \emph{Theorem Proving in Lean 4}, chapter ``Tactics.'' \url{https://lean-lang.org/theorem_proving_in_lean4/Tactics/}
\bibitem{leanaxioms} \emph{The Lean Language Reference}, section ``Axioms.'' \url{https://lean-lang.org/doc/reference/latest/Axioms/}
\bibitem{leanvalidate} \emph{The Lean Language Reference}, section ``Validating a Lean Proof.'' \url{https://lean-lang.org/doc/reference/latest/ValidatingProofs/}
\bibitem{mil} J. Avigad and P. Massot, with contributors, \emph{Mathematics in Lean}. An interactive textbook with exercises, built on the Mathlib library. \url{https://leanprover-community.github.io/mathematics_in_lean/}
\bibitem{pystructures} Python Software Foundation, \emph{The Python Tutorial}, section ``Data Structures.'' \url{https://docs.python.org/3/tutorial/datastructures.html}
\bibitem{pyfractions} Python Software Foundation, \emph{The Python Standard Library}, section ``fractions---Rational numbers.'' \url{https://docs.python.org/3/library/fractions.html}
\end{thebibliography}
\par\medskip\noindent{\small All web addresses were accessed in October 2026. The pages and repository files they point to may change after that date.\par}
